Neural computers point toward systems where a single latent runtime state acts as the computer itself, driving pixels, text, and actions while subsuming what operating systems and interfaces handle today.
Our CNC capability map organizes this view into eight capability axes.
These include efficiency, computation \& reasoning, memory \& storage, I/O \& control, tool bridges, condition-driven generalization, programmability, and artifact generation.
The map is staged and dependency-informed, but not a strict prerequisite chain.
We expect progress in world models and agents to increasingly strengthen and compose these capabilities toward completely neural computers.
In such systems, learned latent dynamics, rather than fixed code and kernels, define the computer’s instruction set, memory, and interfaces.


\section*{Acknowledgements}
The authors also thank Yasheng Sun for his early input of the GUI data collection. 
%
The authors  thank Deyao Zhu and Firas Laakom for their feedback on the manuscript.  
%
Mingchen Zhuge, Haozhe Liu, Shuming Liu, Wenyi Wang, Wenxuan Zhang, Junjie Fei, and Jürgen Schmidhuber were supported by funding from the King Abdullah University of Science and Technology (KAUST) Center of Excellence for Generative AI (award number 5940) and the SDAIA-KAUST Center of Excellence in Data Science and Artificial Intelligence.
